\documentclass[thesis=master, language=german, paper=a4]{isys}[2013/06/28]
%
% [2013/06/28] ensure that the newest isys.cls is used
%              if an older version is used, latex issues a warning
%
% thesis   = sa       - Studienarbeit  / Student thesis 
%            da       - Diplomarbeit   / Diploma thesis
%            bachelor - Bachelorarbeit / Bachelor thesis
%            master   - Masterarbeit   / Master thesis
%            semester - Semesterarbeit / Semester thesis
%
% language = {english, german} 
%            Set default language of document. The class file isys.cls loads the babel
%              package with both options [english, ngerman]. For text that shall be 
%              typeset in a different language than the default language, the 
%              environment 'otherlanguage' or the command '\foreignlanguage' may be used
%              with either of these options (see babel package documentation)
%
% paper    = {a4, letter} - paper size
%
% thesisdraft - mark document as draft without passing "draft" option to 
%                 other packages (e.g. pictures are shown in the document, etc.)
%
% oneside - this option avoids empty pages in the document output 
%                which are necessary for twoside printing, which is useful for 
%                draft versions. 
%                DO NOT USE THIS OPTION WHEN SUBMITTING YOUR THESIS!
%

% mark overful boxes by a 2pt wide rule 
% -- FIX ALL OVERFULL BOXES BEFORE SUBMITTING YOUR THESIS!
\overfullrule=2pt 


%%%%%%%%%%%%%%%%%%%%%%%%%%%%%%%%%%%%%%%%%%%%%%%%%%%%%%%%%%%%%%%%%%%
% Load packages - useful packages, adapt to your own needs
%
% isys.cls loads the "graphicx" package, so \includegraphics can 
%  directly be used to include figures.
%
%%%%%%%%%%%%%%%%%%%%%%%%%%%%%%%%%%%%%%%%%%%%%%%%%%%%%%%%%%%%%%%%%%%
% input encoding - UTF8 or latin1 allow directly typing german umlauts
% make sure the file encoding of all files matches this encoding!
\usepackage[utf8]{inputenc}
% font enconding - do not change this unless you exactly know what you're doing!
\usepackage[T1]{fontenc}

% automatically use correct quotes for the current language: \enquote{text in quotes}
\usepackage{csquotes}
\MakeAutoQuote{»}{«}  % input »text in quotes« is equivalent to \enquote{text in quotes}
\let\q=\enquote 			% use \q{text in quotes}, which is shorter than \enquote{text in quotes}

% better math-commands (e.q.: \dddots, \eqref{})
\usepackage{amsmath,amssymb}                                 
\usepackage{bm}
% better-looking rules in tabulars, \toprule, \midrule, \bottomrule
\usepackage{booktabs}

% use consistent style for units (upright shape)
% \unit[value]{unit}, \unitfrac[value]{unit_numerator}{unit_denominator}
\usepackage{units}

% tables with fixed table width 
\usepackage{tabularx, tabulary}
% improve ragged right alignment with \RaggedRight
\usepackage{ragged2e}
% enable import of eps pictures using pdflatex
%\usepackage{epstopdf}

% enable hyperlinks and bookmarks in PDF file (must typically be loaded as LAST package)
% if used, print hyperlinks in black and suppress boxes around them 
\usepackage[colorlinks=true, allcolors=black, bookmarksnumbered=true]{hyperref}
% to set bookmarks on the first latex run (hyperref needs two),
% provides addtional options for bookmarks (use carefully!)
\usepackage{bookmark}


% Bibliography using classical bibtex
%\bibliographystyle{alphadin}

% better: use biblatex! provides much more options
% use \printbibliography to print the bibliography 
% (search for \printbibliography in this document and uncomment the line, 
%  comment out the line for bibtex)
\usepackage[style=numeric,
						backend=bibtex8,		% or use biber or bibtex?
						bibencoding=latin1,	% .bib file in latin1 or utf8 encoding?
						maxbibnames=99			% give full list of authors (up to 99) in bibliography
						]{biblatex}
\bibliography{references}


%%%%%%%%%%%%%%%%%%%%%%%%%%%%%%%%%%%%%%%%%%%%%%%%%%%%%%%%%%%%%%%%%%%
% Information for title page 
%%%%%%%%%%%%%%%%%%%%%%%%%%%%%%%%%%%%%%%%%%%%%%%%%%%%%%%%%%%%%%%%%%%

% title in the language in which the thesis is written
\title{Schätzverfahren für adaptive Regelung von Hybridgetrieben}
%\subtitle{ }
% title in foreign language (german or english, depending on the language of the thesis)
\foreigntitle{Estimation Procedure for Adaptive Control of Hybrid Gearboxes}
%\foreignsubtitle{Estimation Procedure for Adaptive Control of Hybrid Gearboxes}

\author{Koray Karci}
\major{Maschinenbau}
\date{30.09.2018}		% use either ISO format YYYY-MM-DD or german format DD.MM.YYYY
										% default: current date

\examiner{Prof. Dr.-Ing. Oliver Sawodny}
%\examiner{Prof. Dr.-Ing. Cristina Tar\'{i}n}
%\supervisor{Dipl.-Ing. Supervisor 1}		% give max. 4 \supervisor {...} commands
\supervisor{M.Sc. Michel Bauer}
%\supervisor{Dipl.-Ing. Supervisor 3}
%\supervisor{M.Sc. Supervisor 4}

%% if work on thesis has been conducted with a partner company/institute
%% use \partner and (if desired also \partnerLogo)
% \partnerLogo has no effect as long as no \partner{...} is given
%\partner{Example company name\\
%			Second line company name}
%\partnerLogo{\includegraphics[height=1.4cm, clip, viewport=0 0 30 10]{logos/logoISYS_deutsch}}



%%%%%%%%%%%%%%%%%%%%%%%%%%%%%%%%%%%%%%%%%%%%%%%%%%%%%%%%%%%%%%%%%%%
% Start of document output
%%%%%%%%%%%%%%%%%%%%%%%%%%%%%%%%%%%%%%%%%%%%%%%%%%%%%%%%%%%%%%%%%%%
\begin{document}

% Titelblatt ausgeben / print title page
\pdfbookmark{Titelseite}{title}
\maketitle

%% Englisches Titelblatt in deutschsprachiger Arbeit erstellen
%\begin{otherlanguage}{english}\maketitle \end{otherlanguage}

%% Create german title page in english theses
%\begin{otherlanguage}{ngerman}\maketitle \end{otherlanguage}


	
% Kurzfassung in der Dokumentensprache / Abstract in default document language
\pdfbookmark{Kurzfassung}{abstract-de}
\begin{abstract}
Kurzfassung
\end{abstract}

% Falls gewuenscht, englischen Abstract in deutschsprachiger Arbeit einbinden
\newpage
\begin{otherlanguage}{english}
\pdfbookmark{Abstract}{abstract-en}
\begin{abstract}
Abstract 
\end{abstract}
\end{otherlanguage}

%% Inhalts-/Abbildungs-/Tabellenverzeichnis
%% Table of contents, List of Figures/Tables

	
	
%% If subsections shall also be included in the table of contents
%% uncomment next line. Default: only sections are included.
\setcounter{tocdepth}{1} % include only chapters and sections (1.1) - default
%\setcounter{tocdepth}{2} % also include subsections - use only if told to do so by your advisor!
\cleardoublepage
\pdfbookmark{\contentsname}{toc}
\tableofcontents

\addchap{Abkürzungen und Formelzeichen}

\section*{Abkürzungen}

\begin{flushleft}
\begin{tabularx}{\textwidth}{XlX}	
	&&\\
	FHG								&Freiheitsgrade&\\
	SDG								&Sensitivitätsdifferentialgleichgung&\\
	IMC								&Internal Model Control&\\
	MPC								&Modellprädiktive Regelung (Engl.: Model Predictive Control)&
\end{tabularx}
\end{flushleft}


\vspace{20mm}
\section*{Formelzeichen}


\begin{flushleft}
\begin{tabularx}{\textwidth}{llX}
	\toprule
	Formelzeichen					 	& Einheit									& Beschreibung \\
	\midrule
	$I$											&	A												&	Strom \\
	$U$											&	V												&	Spannung \\
	$T$											&	K												&	Temperatur \\
	&&\\
	$\tau$									&	s												& Zeitkonstante \\
	$\rho_{\mathrm{L}}$			&	$\unitfrac[]{kg}{m^3}$	& Dichte der Luft \\
	&&\\
	$\boldsymbol{u}$						& 												&	Regeleingang des zu regelnden Systems \\
	$\boldsymbol{x}$						& 												&	Zustandsvektor \\
	$\boldsymbol{y}$						& 												&	Regelausgang des zu regelnden Systems \\
\bottomrule
\end{tabularx}
\end{flushleft}


%% Chapters and sections without numbering but included in running heading
%\addchap{Chapter without numbering, included in table of contents}
%\addchap*{Chapter without numbering, not included in table of contents}
%\addsec{Section without numbering, included in table of contents}
%\addsec*{Section without numbering, not included in table of contents}

%%%%%%%%%%%%%%%%%%%%%%%%%%%%%%%%%%%%%%%%%%%%%%%%%%%%%%%%%%%%%%%%%%%
% Beginn des Textes Inhalts / Start of thesis content
%%%%%%%%%%%%%%%%%%%%%%%%%%%%%%%%%%%%%%%%%%%%%%%%%%%%%%%%%%%%%%%%%%%

	
% It may be helpful to put each chapter into a separate file. 
% \includeonly{...} allows to restrict the files that are included,
% which may speed up the document creation process.
% \include always starts a new odd page before including the file 
% contents. If this is not desired, \input can be used.
\chapter{Einleitung und Motivation}\label{ch:ch1}
Dies ist die Einleitung und Motivation.




\chapter{Modellierung}\label{ch:ch2}
In diesem Kapitel wird ein für die Simulation verwendetes detailliertes Modell des Antriebsstrangs hergeleitet. Hierfür wird ein vorhandenes, starres Modell um Elastizitäten der Seitenwellen ergänzt. Für die Herleitung wird dabei der Lagrange Formalismus verwendet, welcher daher auch kurz beschrieben wird. Im weiteren werden aus dem detaillierten Modell vereinfachte Modelle hergeleitet, welche dem Entwurf der Beobachter im Kapitel~\ref{ch:ch4} dienen.


\section{Modellierung des Antriebsstranges}
Zur Modellierung des Antriebsstranges wird zunächst die Getriebedynamik mittels der Newton-Euler Gleichungen berechnet. Dazu wird zu Beginn kurz auf die Herleitung dieser Gleichungen eingegangen. Danach wird dieses um die Elastizitäten der Seitenwellen ergänzt, woraus sich ein schwingfähiges System ergibt.  

\subsection{Herleitung der Newton-Euler Gleichungen}
Die Position und Orientierung von jedem Körper $i$ eines holonomen Mehrkörpersystems mit $p$ Körpern und $f$ Freiheitsgraden, kann in Abhängigkeit der verallgemeinerten Koordinaten \textbf{y} $\in \mathbb{R}^f$ angegeben werden mit
\begin{align}
\pmb{r}_i &= \pmb{r}_i(\pmb{y},t)\quad i = 1(1)p\\
\pmb{s}_i &= \pmb{s}_i(\pmb{y},t).
\end{align}
Durch die zeitliche Ableitung dieser ergeben sich die translatorischen und rotatorischen Geschwindigkeiten
\begin{align}\label{vi}
\pmb{v}_i &=  \dot{\pmb{r}}_i = \frac{\partial\pmb{r}_i}{\partial\pmb{y}}\cdot \dot{\pmb{y}} + \frac{\partial\pmb{r}_i}{\partial\pmb{t}} = \pmb{J}_{Ti}(\pmb{y},t)\cdot \dot{\pmb{y}} + \overline{\pmb{v}}_i(\pmb{y},t)\\
\pmb{\omega}_i &=  \dot{\pmb{s}}_i = \frac{\partial\pmb{s}_i}{\partial\pmb{y}}\cdot \dot{\pmb{y}} + \frac{\partial\pmb{s}_i}{\partial\pmb{t}} = \pmb{J}_{Ri}(\pmb{y},t)\cdot \dot{\pmb{y}} + \overline{\pmb{\omega}}_i(\pmb{y},t)
\end{align}
mit der Jacobi-Matrix der Translation $\pmb{J}_{Ti}$ und der Rotation $\pmb{J}_{Ri}$. Die lokalen Geschwindigkeitsvektoren $\overline{\pmb{v}}_i(\pmb{y},t)$ und $\overline{\pmb{\omega}}_i(\pmb{y},t)$ treten nur bei Systemen mit rheonomen Bindungen auf. Nach erneuter zeitlicher Ableitung ergeben sich die Beschleunigungen 
\begin{align}
\pmb{a}_i = \dot{\pmb{v}}_i &= \pmb{J}_{Ti} \cdot \ddot{\pmb{y}} + \dot{\pmb{J}}_{Ti} \cdot \dot{\pmb{y}} + \dot{\overline{\pmb{v}}}_i = \pmb{J}_{Ti}(\pmb{y},t) \cdot \ddot{\pmb{y}} + \overline{\pmb{a}}_i(\pmb{y},\dot{\pmb{y}},t)\\
\pmb{\alpha}_i = \dot{\pmb{\omega}}_i &= \pmb{J}_{Ri} \cdot \ddot{\pmb{y}} + \dot{\pmb{J}}_{Ri} \cdot \dot{\pmb{y}} + \dot{\overline{\pmb{\omega}}}_i = \pmb{J}_{Ri}(\pmb{y},t) \cdot \ddot{\pmb{y}} + \overline{\pmb{\alpha}}_i(\pmb{y},\dot{\pmb{y}},t)\label{alphai}
\end{align}
mit den lokalen Beschleunigungsvektoren $\overline{\pmb{a}}_i$ und $ \overline{\pmb{\alpha}}_i$. 
Des weiteren lassen sich für die Dynamik eines Systems virtuelle Bewegungen definieren. Diese sind willkürliche, infinitesimale Bewegungen des Systems, welche mit den skleronomen und den rheonomen Bindungen verträglich sind. Für virtuelle Bewegungen an holonomen Bindungen gilt
\begin{align}
\delta \pmb{r} &\neq \pmb{0},\\
\delta \pmb{r} &= \pmb{0},\label{VB0}\\
\delta t &= 0.
\end{align} 
Damit lassen sich die virtuellen Bewegungen der einzelnen Körper $i$ definieren zu
\begin{align}\label{VBr}
\delta \pmb{r}_i = \frac{\partial\pmb{r}_i}{\partial\pmb{y}}\cdot \delta\pmb{y} &= \pmb{J}_{Ti}(\pmb{y},t)\cdot\delta\pmb{y}\\
\delta \pmb{s}_i = \frac{\partial\pmb{s}_i}{\partial\pmb{y}}\cdot \delta\pmb{y} &= \pmb{J}_{Ri}(\pmb{y},t)\cdot\delta\pmb{y}.\label{VBs}
\end{align}
Die Kinetik eines Körpers $i$ wird durch die Newtonschen- und Eulerschen-Gleichungen beschrieben. 
Die Newtonschen-Gleichungen eines Körpers $i$ lassen sich angeben als
\begin{equation}\label{NewtonGl}
m_i \pmb{a}_i(t) = \pmb{f}_i(t)
\end{equation}
wobei $m_i$ die Masse des Körpers angibt und $\pmb{f}_i(t)$ die Summe der angreifenden Kräfte. Diese lassen sich wiederum einteilen in eingeprägte Kräfte $\pmb{f}^e$ und Reaktionskräfte $\pmb{f}^r$.
Die Eulerschen Gleichungen lauten
\begin{equation}\label{EulerGl}
\pmb{I}_i(t)\cdot\pmb{\alpha}_i(t) + \tilde{\pmb{\omega}}_i(t)\cdot \pmb{I}_i(t)\pmb{\omega}_i(t)= \pmb{l}_i(t)
\end{equation}
mit dem Trägheitstensor $\pmb{I}_i(t)$ und den äußeren Momenten $\pmb{l}_i(t)$ bezogen auf den Massenmittelpunkt des Körpers $i$ im Inertialsystem und dem schiefsymmetrische 3x3-Tensor $\tilde{\pmb{\omega}}_i(t)$. Die äußeren Momenten $\pmb{l}_i(t)$ lassen sich äquivalent zum translatorischen Fall wieder in eingeprägte Momente $\pmb{l}^e$ und Reaktionsmomente $\pmb{l}^r$ aufteilen. Der schiefsymmetrische 3x3-Tensor ist für ein beliebigen Vektor 
\begin{equation}
\pmb{v} = [v_1 \quad v_2 \quad v_3]^T
\end{equation}
definiert als
\begin{equation}
\tilde{\pmb{v}} = \begin{bmatrix} 0 & -v_3 & v_2 \\ v_3 & 0 & -v_1 \\ -v_2 & v_1 & 0 \end{bmatrix}.
\end{equation}
Darüber hinaus kann mit den virtuellen Bewegungen aus (\ref{VBr}) und (\ref{VBs}) die virtuelle Arbeit $\delta \pmb{W}$ definiert werden. Da es, wie in (\ref{VB0}) definiert, in Richtung der Bindungen keine virtuellen Bewegungen geben kann, ergibt sich für die virtuelle Arbeit der Reaktionskräfte
\begin{equation}\label{VA}
\delta\pmb{W}^r = \sum_{i=1}^p(\pmb{f}_i^r\cdot \delta \pmb{r}_i + \pmb{l}^r_i\cdot\delta\pmb{s}_i)=0.
\end{equation}
Mit den Gleichungen \ref{NewtonGl}, \ref{EulerGl} und der beschriebenen Aufteilung der äußeren Kräfte in eingeprägte Kräfte und Reaktionskräfte, folgt aus (\ref{VA}) das d’Alembertsche Prinzip für Mehrkörpersysteme
\begin{equation}\label{VA}
\sum_{i=1}^p\left[ (m_i\pmb{a}_i-\pmb{f}_i^e)\cdot \delta \pmb{r}_i + (\pmb{I_i}\cdot \pmb{\alpha}_i + \tilde{\pmb{\omega}}\cdot\pmb{I}_i\cdot\pmb{\omega}_i-\pmb{l}^e_i)\cdot\delta\pmb{s}_i\right] =0.
\end{equation}
Durch einsetzten der virtuellen Bewegungen (\ref{VBr}) und (\ref{VBs}) ergibt sich daraus
\begin{equation}
\delta\pmb{y}\sum_{i=1}^p\left[\pmb{J}^T_{Ti}\cdot(m_i\pmb{a}_i-\pmb{f}_i^e)+ \pmb{J}^T_{Ri}\cdot(\pmb{I_i}\cdot \pmb{\alpha}_i + \tilde{\pmb{\omega}}\cdot\pmb{I}_i\cdot\pmb{\omega}_i-\pmb{l}^e_i)\right]=0.
\end{equation}
Das einarbeiten der kinematischen Gleichungen (\ref{vi}) -- (\ref{alphai}) in ... führt zur Bewegungsgleichung für holonome Mehrkörpersysteme (MKS) 
\begin{align}\label{Bwg_lang}
\begin{split}
\sum_{i=1}^p\left[\pmb{J}_{Ti}^T\cdot m_i \cdot \pmb{J}_{Ti} + \pmb{J}_{Ri}^T\cdot \pmb{I}_i \cdot \pmb{J}_{Ri}\right]\ddot{\pmb{y}}+\sum_{i=1}^p\left[\pmb{J}_{Ti}^T\cdot m_i \cdot \overline{\pmb{a}}_i + \pmb{J}^T_{Ri}\cdot\pmb{I}_i\cdot\overline{\pmb{\alpha}}_i+\pmb{J}^T_{Ri}\cdot\tilde{\pmb{\omega}}_i\cdot\pmb{I}_i\cdot\pmb{\omega_i}\right]=\\ \sum_{i=1}^p\left[\pmb{J}^T_{Ti}\cdot\pmb{f}_i^e + \pmb{J}_{Ri}^T\cdot\pmb{I}_i^e\right].
\end{split}
\end{align}
Darin lässt sich die erste Summe zur Massenmatrix $\pmb{M}(\pmb{y},t)$, die zweite Summe zum Vektor der verallgemeinerten Zentrifugal-- und Coriloriskräfte $\pmb{k}(\pmb{y},\dot{\pmb{y}},t)$ und die dritte Summe zum Vektor der verallgemeinerten Kräfte $\pmb{q}(\pmb{y},\dot{\pmb{y}},t)$ zusammenfassen, wodurch sich (\ref{Bwg_lang}) schreiben lässt als
\begin{equation}\label{Bwg}
\pmb{M}(\pmb{y},t)\cdot\ddot{\pmb{y}}+\pmb{k}(\pmb{y},\dot{\pmb{y}},t) = \pmb{q}(\pmb{y},\dot{\pmb{y}},t).
\end{equation}
  

\subsection{Berechnung der Getriebedynamik}
Mit der allgemeinen Bewegungsgleichung für MKS (\ref{Bwg}) kann die Getriebedynamik berechnet werde. Das in dieser Arbeit betrachtete Getriebe besteht im wesentlichen aus vier hintereinander liegenden Umlaufgetrieben. Der schematische Aufbau eines Umlaufgetriebes ist in Abbildung ??? zu sehen. Dieses besteht aus dem Sonnenrad in der Mitte, den darum gleichmäßig verteilten Planetenrädern, welche auf dem Planetenträger drehbar gelagert sind und dem Hohlrad. Der Zusammenhang der einzelnen Drehzahlen ist gegeben durch die beiden Zwangsbedingungen
\begin{align}
\mathrm{Sonne - Planeten}:\quad &r_S\omega_S - r_T\omega_T+ r_P\omega_P=0 \\
\mathrm{Planeten - Hohlrad}:\quad &r_H\omega_H - r_T\omega_T- r_P\omega_P=0.
\end{align}
Dabei gilt, dass die Verhältnisse des der Rollradien denen der Zähnezahlen einer Zahnradverbindung entsprechen. Der Rollradius des Planetenträger entspricht somit der Summe der Zähnezahl der Sonne $z_S$ und des Planeten $z_P$ und es ergibt sich
\begin{align}
\mathrm{Sonne - Planeten}:\quad &z_S\omega_S - (z_S + z_P)\omega_T+ z_P\omega_P=0 \\
\mathrm{Planeten - Hohlrad}:\quad &z_H\omega_H - (z_S + z_P)\omega_T- z_P\omega_P=0.
\end{align}

\subsection{Erweiterung um elastische Seitenwellen}

\section{Modellierung der Reibung an Kupplungen und Bremsen}
\chapter{Sensitivitätsanalyse}\label{ch:Sensitivity}\label{ch:ch3}
In diesem Kapitel wird das in Kapitel \cite{ch:ch3} hergeleitete lineare Modell des Antriebsstrangs auf die Sensitivität verschiedener Parameter und Eingänge untersucht. Damit soll die Relevanz für eine Parameter- und Eingangsschätzung im folgenden Kapitel belegt und unterstrichen werden.
Eine mögliche Definition für eine Sensitivitätsanalyse ist die folgende: \emph{The study of how uncertainty in the output of a model (numerical or otherwise) can be apportioned to different sources of uncertainty in the model input} \cite{Saltelli.2004}.
Für die Analyse gibt es verschiedene Methoden. Dabei ist die OAT(one-at-a-time)-Methode wohl die einfachste. Dabei wird pro Simulationsdurchlauf lediglich eine Eingangsgröße geändert während alle anderen auf den nominal Werten gehalten werden. Die Methode liefert aber nur eine sehr beschränkte Aussage über den gesamten Eingangsraum. Jedoch kann sie für die Abschätzung einer ersten Untergruppe an relevanten Eingängen hilfreich sein.    

 Die in der Literatur am häufigsten diskutierten Ansätze basieren auf partiellen Ableitungen der Art $\partial y_i/\partial x_i$ wobei $y$ ein beliebiger Ausgang und $x_i$ ein beliebiger Eingang des Systems ist. Ein Eingang kann hierbei auch ein variierender Parameter sein. Dieses Vorgehen wird vor allem bei relativ einfachen, mathematisch wohlbestimmten und stetigen Modellen verwendet. Der Vorteil dabei ist, dass die Sensitvitätsfunktionen sobald sie einmal berechnet sind leicht an Veränderungen anpassbar sind und auch für ähnliche Systeme verwendet werden können \cite{Karnavas.1993}. Jedoch werden die partiellen Ableitung immer in einem bestimmen Arbeitspunkt berechnet d.h. die Aussagen der Sensitvitätsfunktionen gelten für nichtlineare Systeme nur lokal um diesen den Arbeitspunkt. Des weiteren ist es oft sehr schwer oder gar unmöglich die benötigten Ableitungen für komplexe Modelle zu berechnen.

Eine Alternative zur Berechnung der partiellen Ableitungen sind globale Methoden. Hierbei werden die Modelleingänge stochstisch varriert, woraus sich nach vielen Simulationen eine Strichprobe für die Modellausgänge ergibt. Dieses Methode der mehrfachen Simulation wird auch Monte-Carlo-Methode genannt. Wie der Name schon sagt ist der grundlegende Vorteil dabei, dass durch eine genügend große Stichprobe eine globale Aussage zur Sensitivität gegeben werden kann. Des weiteren erlaubt die Methode auch Aussagen über die verkoppelte Wirkung von Eingangsänderungen ohne spezielles Vorwissen über die Struktur des Modells. Der entscheidende Nachteil liegt je nach Stichprobengröße und Komplexität des Modells in der Rechenzeit.      

Aufgrund des linearen und stetigen Charakters von \eqref{eq:sys_linex} werden für diese Arbeit die partiellen Ableitungen für die Sensitivitätsanalyse verwendet.
Hierfür wird zuerst ein idealer Schaltvorgang beschrieben um die Sensitivitätsfunktionen in verschieden Schaltphasen auswerten zukönnen. Des weiteren wird mit der Berechnung der Fisher-Informationsmatrix eine Abschätzung für die erreichbare Genauigkeit der Schätzungen geliefert.
   
\section{Beschreibung eines Gangwechsels}\label{sec:Gangwechsel}
Bei einem Gangwechsel wird durch die Betätigung von Schaltelementen die Übersetzung des Getriebes geändert. Dies soll zum einen in jedem Fall möglichst ohne für den Fahrer spürbare unerwünschte Beschleunigungen am Getriebeabtrieb passieren. Zum anderen soll der Gangwechsel aber auch so schnell wie möglich und ohne Momentenunterbrechung vollzogen werden. Es besteht also eine Zielkonflikt zwischen den Komfortansprüchen des Fahrers und der Dauer eines Schaltvorgangs. Grundsätzlich wird zwischen vier verschieden Schaltvorgängen unterschieden. 
Bei dem in dieser Arbeit betrachteten Getriebe geschieht ein Gangwechsel mittels der Schließung und der simultanen Öffnung einer Bremse bzw. Kupplung. In der   

\section{Anwendung der Paramater- und Eingangssensitivitätsanalyse}
Wie einleitend erläutert werden die Sensitiviäten in dieser Arbeit mittels partiellen Differentialen berechnet. Nach \cite{Hearne.1985} ist für ein nichtlineares System erster Ordnung
\begin{equation}
\mathrm{d} x_i/\mathrm{d} t = f_i(\pmb{x},\pmb{p},\pmb{u})\quad i=1,2,\,\dots\, ,n 
\end{equation}
mit dem Zustandsvektor $\pmb{x} = [x_1, x_2,\, \dots\, ,x_n]^T$, dem Parametervektor $\pmb{p} = [p_1, p_2,\, \dots\, ,p_m]^T$, dem Eingangsvektor $\pmb{u} = [u_1, u_2,\, \dots\, ,p_l]^T$ und der differenzierbaren Funktion $f_i$ die absolute Parametersensitivitätsfunktion definiert als
\begin{equation}\label{eq:abs_psens}
S_{ij} = \partial x_i/\partial p_j \quad i=1,2,\,\dots\, ,n \quad j=1,2,\,\dots\, ,m.
\end{equation}
Die entsprechende normalisierte Parametersensitivitätsfunktion ist 
\begin{equation}
S_{\mathrm{N},ij} = \partial x_i/\partial p_j\ p_{j0}/x_{i0}.
\end{equation}
Die normalisierte Parametersensitivitätsfunktion $S_{ij}$ kann interpretiert werden, als approximierte prozentuale Änderung des Zustandes $x_i$ für eine $1\%$-ige Änderung des Parameters $p_j$. Dem entsprechend sind auch die absolute Eingangssensitivitätsfunktion
\begin{equation}\label{eq:abs_usens}
S_{ik} = \partial x_i/\partial u_k \quad i=1,2,\,\dots\, ,n \quad k=1,2,\,\dots\, ,l
\end{equation}
und die normalisierte Eingangssensitivitätsfunktion
\begin{equation}\label{eq:abs_usens}
S_{ik} = \partial x_i/\partial u_k\ u_{k0}/x_{i0}
\end{equation}
definiert.
Wendet man \eqref{eq:abs_psens} auf das lineare System 
\begin{equation}
\mathrm{d} \pmb{x}/\mathrm{d} t = \pmb{A}(\pmb{p})\,\pmb{x}+\pmb{B}(\pmb{p})\,\pmb{u}
\end{equation}
an, erhält man
\begin{equation}
\mathrm{d} \pmb{S}_j/\mathrm{d} t = \dot{\pmb{S}}_j= \pmb{A}(\pmb{p})\,\pmb{S}_j+\partial/\partial p_j\ \pmb{B}(\pmb{p})\,\pmb{u}.
\end{equation}
Für die Eingangssensitivität ergibt sich 
\begin{equation}
\mathrm{d} \pmb{S}_k/\mathrm{d} t = \dot{\pmb{S}}_k= \partial/\partial u_i\ \pmb{B}(\pmb{p})\,\pmb{u}.
\end{equation}
Die Parameter- und Eingangsensitivitätsfunktion werden im Folgenden anhand \eqref{eq:sys_ex23lin} für $m_\mathrm{veh}$, $k_\mathrm{ss}$, $d_\mathrm{ss}$, $\mu_{38}$ und $\zeta$ gebildet und die zeitlichen Verläufe von $\pmb{S}$ für einen Schaltvorgang gemäß \ref{sec:Gangwechsel} interpretiert.

\subsection{Sensitivität des Zustandes $\pmb{x}$ von $k_\mathrm{ss}$}

kann die Sensitivitätsfunktion wie fo 


\section{Fisher-Informationsmatrix}





\chapter{Beobachterentwurf}\label{ch:ch3}


\section{Kalman-Filter}
\section{Extented Kalman-Filter}
\section{Unknown Input Observer}
\section{Extended Unknown Input Observers}





\chapter{Zusammenfassung und Ausblick}\label{ch:ch4}
Im nächsten Schritt werden zwei verschiedene Online-Parameterschätzverfahren beschrieben und die dazugehörigen Simulationsergebnisse diskutiert. Die Parameter werden dabei in den verschieden Schaltphasen geschätzt.

Des weiteren kann durch die Betrachtung der Widerstandskräfte $F_\mathrm{r}(v)$ und $F_\mathrm{ad}(v)$ als Eingangsgrößen, das System weiterhin in lineare Form
\begin{equation}\label{eq:sys_linex}
\dot{\pmb{x}}_\mathrm{ex} = \pmb{A}_\mathrm{ex}\,\pmb{x}_\mathrm{ex} + \pmb{B}_\mathrm{ex}\,\pmb{u}_\mathrm{ex}
\end{equation}
angegeben werden. Hierfür werden die beiden Widerstandskräfte in \eqref{eq:dynwc} durch $T^*_\mathrm{res}$ ersetzt, sodass sich 
\begin{equation}\label{eq:dynwc_lin}
\dot{\omega}_\mathrm{C} = \left[ 2\,\left(k_\mathrm{ss}\left(\frac{\phi_2}{i_{\mathrm{D}}} - \phi_\mathrm{C}\right) + d_\mathrm{ss}\left(\frac{\omega_2}{i_{\mathrm{D}}} - \omega_\mathrm{C}\right)\right) - T^*_\mathrm{res} - r_\mathrm{dyn}\, F_\mathrm{g}\right]/I^\mathrm{eff}_\mathrm{C}.
\end{equation}
 ergibt. Der Eingangsvektor des linearen Systems ist dann definiert als 
\begin{equation}
\pmb{u}_\mathrm{ex} = \begin{bmatrix} \mathrm{T}_{\mathrm{In}} & T^*_\mathrm{res} & \mathrm{T}_{\mathrm{K81}} & \mathrm{T}_{\mathrm{K38}} & \mathrm{T}_{\mathrm{B08}} & \mathrm{T}_{\mathrm{B05}} & \mathrm{T}_{\mathrm{B06}} & \mathrm{T}_{\mathrm{K27}} & \zeta \end{bmatrix}^T.
\end{equation}

Die Funktion $\pmb{f}(\pmb{x}_\mathrm{ex})$ und die Eingangmatrix $\pmb{B}_\mathrm{ex,nl}$ des nichtlinearen Systems, sowie die Systemmatrix $\pmb{A}_\mathrm{ex}$ und die Eingangmatrix $\pmb{B}_\mathrm{ex}$ des linearen Systems sind im Anhang mit numerischen Werten für die Zähnezahlen der Zahnräder und der Massenträgheitsmomente angegeben.


Mit den eingesetzten numerischen Werten der Zahnräder, der Massenträgheitsmomente, der Übersetzung der Differentials und dem dynamischen Radius der Räder ergeben sich die Systemmatrizen zu
\begin{align}
\pmb{A}_\mathrm{ex,23} &= \begin{bmatrix} 0 & -0,799\,d_\mathrm{ss} & -1,97\,k_\mathrm{ss} & 1,97\,d_\mathrm{ss} \\
0 & -0,248\,d_\mathrm{ss} & -0,612\,k_\mathrm{ss} & 0,612\,d_\mathrm{ss}\\
0 & 0,405 & 0 & -1,0\\
0 & \frac{0,81\,d_\mathrm{ss}}{m_\mathrm{veh}\,0,108 + 6,4} & \frac{2,0\,k_\mathrm{ss}}{m_\mathrm{veh}\,0,108 + 6,4} & \frac{-2,0\,d_\mathrm{ss}}{m_\mathrm{veh}\,0,108 + 6,4})\end{bmatrix}\\
\nonumber \\
\pmb{B}_\mathrm{ex,23} &= \begin{bmatrix} 10,7 & -2,79 & -17,0 & 0 & 0\\
2,44 & 0,0135 & -2,41 & 0 & 0 \\
0 & 0 & 0 & 0 & 0\\
0 & 0 & 0 & \frac{-1}{0,108\,m_\mathrm{veh} + 6,4} & \frac{-3,22\,m_\mathrm{veh}}{0,108\,m_\mathrm{veh} + 6,4} \end{bmatrix}
\end{align}
\cleardoublepage

\appendix
\chapter{Anhang}\label{ch:appCh1}
Im Anhang werden Ergebnisse aufgeführt, die zwar im Kontext der Arbeit wesentlich, jedoch nicht für das Verständnis der Arbeit notwendig sind. Beispielsweise sind dies detailliertere Messergebnisse, Herleitungen oder Datenblätter. 

% List of tables and list of figures at end
\listoffigures
\listoftables

% Literaturverzeichnis
% classical bibtex
%\bibliography{references}

% modern biblatex
\printbibliography[heading=bibintoc]

% in any case ...
\label{ch:bib}



% Erklaerung des Autors, netig fuer Bachelor-/Master-Arbeiten der Universitaet Stuttgart
% Declaration of authorship, necessary for bachelor and master thesis at University of Stuttgart
\DeclarationOfAuthorship
\end{document}
